\documentclass[10pt]{scrartcl}
\usepackage[mdthm,colorsec,noauthor,noasy]{evan}
\usepackage{handout}
\newcommand{\aug}[2]{\left[\begin{array}{@{}#1@{}}#2\end{array}\right]}
\newcommand{\vv}[1]{\mathbf{#1}}
\DeclareMathOperator{\spn}{Span}\DeclareMathOperator{\Nul}{Nul}\DeclareMathOperator{\Col}{Col}\DeclareMathOperator{\Row}{Row}
\title{Math 54 --- Week 6}
\subtitle{Linearly independent sets and bases; dimension and rank (Lay 4.3, 4.5)}
\author{}\date{}\ihead{}
\begin{document}
\maketitle

\section{Facts used from Weeks 1--5}
\begin{enumerate}[label=(W\arabic*),leftmargin=2.8em,itemsep=4pt]
  \ii In a dependent set $\{\vv v_1,\dots,\vv v_p\}$ with $\vv v_1\ne\vv0$, some $\vv v_j$ is a combination of $\vv v_1,\dots,\vv v_{j-1}$. More vectors than entries in $\RR^n$ are dependent.
  \ii Columns of $A$ independent $\iff$ pivot in every column. $A\vv x=\vv0$ has a nontrivial solution $\iff$ a free variable. A homogeneous system with more unknowns than equations has a nontrivial solution.
  \ii Row operations do not change the solution set of $A\vv x=\vv0$.
  \ii Subspaces; $\spn$ is a subspace; $\Nul A$, $\Col A$, $\Row A$. The parametric vector form of $A\vv x=\vv0$ gives a spanning set for $\Nul A$, one vector per free variable, and that set is independent.
\end{enumerate}

\section{This week}

\subsection{Linear independence in a vector space; bases (Lay 4.3)}
\begin{itemize}
  \ii The definitions carry over unchanged: $\{\vv v_1,\dots,\vv v_p\}$ in $V$ is \vocab{linearly independent} if $c_1\vv v_1+\dots+c_p\vv v_p=\vv0$ only for $c_1=\dots=c_p=0$. The theorem (W1) about some $\vv v_j$ being a combination of the earlier vectors holds in any vector space, with the same proof.
  \ii In $\mathbb P_2$, $\{1,\ t,\ t^2\}$ is independent: $c_0+c_1t+c_2t^2$ is the zero polynomial only if all coefficients are $0$. $\{\sin t\cos t,\ \sin2t\}$ is dependent in the space of functions, since $\sin2t=2\sin t\cos t$.
\end{itemize}
\begin{qbox}
  A spanning set may contain redundant vectors. What is the smallest kind of set that still spans, and how do we find one for $\Nul A$ and $\Col A$?
\end{qbox}
\begin{definition*}
  A \vocab{basis} for a subspace $H$ of $V$ is a set $\mathcal B=\{\vv b_1,\dots,\vv b_p\}$ in $H$ that is linearly independent and spans $H$.
\end{definition*}
\begin{itemize}
  \ii The columns $\vv e_1,\dots,\vv e_n$ of $I_n$ form the \vocab{standard basis} of $\RR^n$. $\{1,t,\dots,t^n\}$ is the standard basis of $\mathbb P_n$.
  \ii The columns of an invertible $n\times n$ matrix form a basis of $\RR^n$: by the invertible matrix theorem they are independent and span $\RR^n$.
\end{itemize}
\begin{theorem*}[Spanning set theorem]
  Let $S=\{\vv v_1,\dots,\vv v_p\}$ and $H=\spn S$. (a) If some $\vv v_k$ is a combination of the others, removing it from $S$ does not change the span. (b) If $H\ne\{\vv0\}$, some subset of $S$ is a basis for $H$.
\end{theorem*}
\begin{proof}
  \begin{steps}
    \ii (a) Write $\vv v_k=\sum_{i\ne k}a_i\vv v_i$. In any combination $\sum c_i\vv v_i$, replace $\vv v_k$ by this expression; the result uses only the other vectors.
    \ii (b) If $S$ is independent it is a basis. Otherwise some vector is a combination of the others (W1); remove it by (a). Repeat. The process stops, at the latest when one nonzero vector remains, which is independent.
  \end{steps}
\end{proof}
A basis is a spanning set with nothing removable, and an independent set to which nothing in $H$ can be added without losing independence.

\textbf{A basis for $\Nul A$}: the vectors from the parametric vector form of $A\vv x=\vv0$ (W4).

\begin{theorem*}[A basis for $\Col A$]
  The pivot columns of $A$ form a basis for $\Col A$.
\end{theorem*}
\begin{ideabox}
Row reduction changes the columns, so $\Col A$ and $\Col U$ are usually different spaces. But it does not change \emph{which combinations of the columns are zero}: a relation $c_1\vv a_1+\dots+c_n\vv a_n=\vv0$ is a solution of $A\vv x=\vv0$, and $A$ and its echelon form $U$ have the same solutions (W3). In $U$ it is visible that the pivot columns are independent and every other column is a combination of them. The same relations therefore hold among the columns of $A$.
\end{ideabox}
\begin{proof}
  \begin{steps}
    \ii Let $B$ be the reduced echelon form of $A$. $A\vv x=\vv0$ and $B\vv x=\vv0$ have the same solutions, so a list of weights gives $\vv0$ on the columns of $A$ exactly when it gives $\vv0$ on the same columns of $B$.
    \ii The pivot columns of $B$ are distinct standard vectors $\vv e_i$, so they are independent; each non-pivot column of $B$ is a combination of the pivot columns before it.
    \ii By step 1, the pivot columns of $A$ are independent, and each non-pivot column of $A$ is the same combination of pivot columns of $A$. By the spanning set theorem the non-pivot columns can be removed without changing $\Col A$.
  \end{steps}
\end{proof}
\textbf{Worked example.}
\[ A=\begin{bmatrix}1&4&0&2&-1\\3&12&1&5&5\\2&8&1&3&2\\5&20&2&8&8\end{bmatrix}\ \longrightarrow\ \text{RREF }\begin{bmatrix}1&4&0&2&0\\0&0&1&-1&0\\0&0&0&0&1\\0&0&0&0&0\end{bmatrix}. \]
\begin{steps}
  \ii Pivot columns 1, 3, 5. Basis for $\Col A$: $\vv a_1=(1,3,2,5)$, $\vv a_3=(0,1,1,2)$, $\vv a_5=(-1,5,2,8)$ (columns of $A$, not of the echelon form).
  \ii $x_2,x_4$ free; $x_1=-4x_2-2x_4$, $x_3=x_4$, $x_5=0$. Basis for $\Nul A$: $(-4,1,0,0,0)$, $(-2,0,1,1,0)$.
  \ii Basis for $\Row A$: the nonzero rows of an echelon form, here $(1,4,0,2,0)$, $(0,0,1,-1,0)$, $(0,0,0,0,1)$. (Row operations replace rows by combinations of rows and can be undone, so they do not change the row space; the nonzero rows of an echelon form are independent, since each has a leading entry in a column where the rows below have zeros.)
\end{steps}

\subsection{Dimension (Lay 4.5)}
\begin{qbox}
  A subspace has many bases. Do they all have the same number of vectors? If so, that number measures the size of the subspace.
\end{qbox}
\begin{theorem*}[More vectors than a basis are dependent]
  If $V$ has a basis $\{\vv b_1,\dots,\vv b_n\}$, then any set of more than $n$ vectors in $V$ is linearly dependent.
\end{theorem*}
\begin{ideabox}
Write each of the vectors $\vv u_1,\dots,\vv u_p$ ($p>n$) in terms of the basis. Finding a dependence relation among the $\vv u$'s then amounts to solving a homogeneous linear system with $n$ equations (one per basis vector) and $p$ unknowns (one weight per $\vv u_j$). More unknowns than equations always gives a nontrivial solution (W2).
\end{ideabox}
\begin{proof}
  \begin{steps}
    \ii Write $\vv u_j=a_{1j}\vv b_1+\dots+a_{nj}\vv b_n$ for $j=1,\dots,p$, and let $M=(a_{ij})$, an $n\times p$ matrix.
    \ii Since $p>n$, $M\vv c=\vv0$ has a nontrivial solution $\vv c=(c_1,\dots,c_p)$.
    \ii Then $\sum_jc_j\vv u_j=\sum_jc_j\sum_ia_{ij}\vv b_i=\sum_i\bigl(\sum_ja_{ij}c_j\bigr)\vv b_i=\sum_i(M\vv c)_i\vv b_i=\vv0$, with the $c_j$ not all $0$.
  \end{steps}
\end{proof}
\begin{theorem*}[All bases have the same size]
  If $V$ has a basis of $n$ vectors, every basis of $V$ has exactly $n$ vectors.
\end{theorem*}
\begin{proof}
  If $\mathcal C$ is another basis, it is independent, so it has at most $n$ vectors by the previous theorem. Exchanging the roles of the two bases, $n$ is at most the size of $\mathcal C$.
\end{proof}
\begin{definition*}
  If $V$ is spanned by a finite set, $V$ is \vocab{finite-dimensional}, and $\dim V$ is the number of vectors in any basis. $\dim\{\vv0\}=0$. Otherwise $V$ is \vocab{infinite-dimensional}.
\end{definition*}
\begin{itemize}
  \ii $\dim\RR^n=n$; $\dim\mathbb P_n=n+1$; the space of all polynomials is infinite-dimensional.
  \ii Subspaces of $\RR^3$ by dimension: $0$: $\{\vv0\}$; $1$: lines through $\vv0$; $2$: planes through $\vv0$; $3$: $\RR^3$ itself.
\end{itemize}
\begin{theorem*}[Subspaces; extending to a basis]
  If $H$ is a subspace of a finite-dimensional $V$, any independent set in $H$ can be extended to a basis of $H$, and $\dim H\le\dim V$.
\end{theorem*}
\begin{proof}
  If an independent set $S$ in $H$ does not span $H$, pick $\vv u\in H$ outside $\spn S$; then $S\cup\{\vv u\}$ is still independent (in a relation, the weight on $\vv u$ must be $0$, else $\vv u\in\spn S$; then all weights are $0$). Repeat. By the theorem above the size never exceeds $\dim V$, so the process stops, with a basis of $H$.
\end{proof}
\begin{theorem*}[Basis theorem]
  Let $V$ have dimension $p\ge1$. Any $p$ linearly independent vectors in $V$ form a basis. Any $p$ vectors that span $V$ form a basis.
\end{theorem*}
\begin{proof}
  An independent set of $p$ vectors extends to a basis, which has exactly $p$ vectors, so nothing was added. A spanning set of $p$ vectors contains a basis (spanning set theorem), which has exactly $p$ vectors, so nothing was removed.
\end{proof}
So when you already know the dimension, you only need to check \emph{one} of the two conditions.

\subsection{Rank (Lay 4.5)}
\begin{itemize}
  \ii $\dim\Nul A$ = number of free variables. $\dim\Col A$ = number of pivot columns. $\dim\Row A$ = number of pivots (the nonzero rows of an echelon form).
  \ii The \vocab{rank} of $A$ is $\rank A=\dim\Col A$. By the line above, $\rank A=\dim\Row A=\rank A^T$.
\end{itemize}
\begin{theorem*}[Rank theorem]
  For an $m\times n$ matrix $A$: $\rank A+\dim\Nul A=n$.
\end{theorem*}
\begin{proof}
  Every column of $A$ is either a pivot column (counted by $\rank A$) or corresponds to a free variable (counted by $\dim\Nul A$). There are $n$ columns.
\end{proof}
\begin{itemize}
  \ii Example: a $7\times9$ matrix with a two-dimensional null space has rank $7$, so $\Col A$ is a 7-dimensional subspace of $\RR^7$, i.e.\ all of $\RR^7$.
  \ii A $6\times9$ matrix has rank at most $6$, so its null space has dimension at least $3$.
\end{itemize}
\textbf{More statements for the invertible matrix theorem} ($A$ is $n\times n$): (m) the columns of $A$ form a basis of $\RR^n$; (n) $\Col A=\RR^n$; (o) $\rank A=n$; (p) $\dim\Nul A=0$; (q) $\Nul A=\{\vv0\}$.

\begin{shortcuts}
  \ii \textbf{Bases for the three spaces from one row reduction}: $\Col A$: pivot columns of $A$; $\Row A$: nonzero rows of the echelon form; $\Nul A$: vectors from the parametric form.
  \ii \textbf{Counting}: $\rank A=\#$pivots $=\dim\Col A=\dim\Row A$; $\dim\Nul A=n-\rank A$.
  \ii \textbf{Is this set a basis of $\RR^n$?} With exactly $n$ vectors, check independence only (or $\det\ne0$). With a different number of vectors: no.
  \ii \textbf{Basis of a subspace given as a span}: put the vectors as columns and keep the pivot columns.
  \ii \textbf{Extending to a basis of $\RR^n$}: row reduce $[\,\vv v_1\cdots\vv v_k\ \vv e_1\cdots\vv e_n\,]$ and keep the pivot columns; the $\vv v$'s come first, so they are all kept.
  \ii \textbf{Polynomial subspaces}: write the conditions as equations on the coefficients; solve; read off a basis.
  \ii \textbf{Rank questions}: $\rank A\le\min(m,n)$. Use the rank theorem before computing anything.
\end{shortcuts}
\begin{trapbox}
  \ii The basis for $\Col A$ is the pivot columns \emph{of $A$}. The columns of the echelon form span a different space in general.
  \ii The basis for $\Row A$ is rows \emph{of the echelon form}; the corresponding rows of $A$ need not work (row swaps change which rows they are).
  \ii A basis is not unique; its size is.
  \ii $\dim\RR^2=2$, but $\RR^2$ is not a subspace of $\RR^3$; a plane through $\vv0$ in $\RR^3$ is a 2-dimensional subspace.
  \ii The rank theorem has $n$ (columns), not $m$.
\end{trapbox}

\section{Homework this week}
HW5, Section 4.3: 3, 5, 13, 15, 20, 23, 24, 26, 29, 31, 41, 44. HW6, Section 4.5: 1, 3, 9, 11, 17, 18, 22, 24, 25, 27, 33, 35, 44, 45 (all the true/false questions in 4.5 are recommended). (Section 4.4 of HW6 is in the Week 7 handout.) Midterm 1 is this week.

\section{Problem set}
\begin{prob}\textbf{P1.}\ For $A=\begin{bmatrix}-2&4&-2&-4\\2&-6&-3&1\\-3&8&2&-3\end{bmatrix}$, find bases for $\Col A$, $\Row A$ and $\Nul A$, and check the rank theorem.
\end{prob}
\begin{prob}\textbf{P2.}\ Is $\{(1,0,-1),(2,1,0),(1,3,4)\}$ a basis of $\RR^3$?
\end{prob}
\begin{prob}\textbf{P3.}\ Find a basis for, and the dimension of, $H=\spn\{(1,0,2),(2,1,5),(0,-1,-1)\}$.
\end{prob}
\begin{prob}\textbf{P4.}\ Find a basis for $\{p\in\mathbb P_3: p(1)=0\}$ and its dimension.
\end{prob}
\begin{prob}\textbf{P5.}\ (a) $A$ is $5\times8$ with rank $5$. What is $\dim\Nul A$? Is $\Col A=\RR^5$? (b) Can a $4\times6$ matrix have a one-dimensional null space?
\end{prob}
\begin{prob}\textbf{P6.}\ Extend $\{(1,2,0),(0,1,3)\}$ to a basis of $\RR^3$.
\end{prob}
\begin{prob}\textbf{P7.}\ True or false, with a reason. (a) If $\dim V=3$, any 4 vectors in $V$ are dependent. (b) If $\dim V=3$, any 2 vectors in $V$ span a plane. (c) The number of pivot columns of $A$ equals $\dim\Row A$. (d) If $A$ is $3\times3$ and $\dim\Nul A=1$, then $\det A=0$.
\end{prob}

\end{document}
